\chapter{结论与展望}

\section{主要结论}
围绕 FAST 主动反射面形状调节问题，本文构建了“理想抛物面确定 -- 节点约束优化 -- 接收比评估”的统一模型框架。主要结论包括：
\begin{enumerate}
  \item 题目可抽象为大规模几何约束优化问题，目标函数与约束均具有明确物理意义；
  \item 通过节点离散与促动器约束耦合，可输出题目要求的节点坐标与伸缩量结果；
  \item 通过接收比指标可对调节方案进行量化验证，实现“结构可行 + 光学有效”的双重评价。
\end{enumerate}

\section{工程意义}
该模型为 FAST 类主动反射面系统提供了可解释、可复算的调节分析路径：既可用于竞赛与课程报告，也可作为工程原型的算法基础模块。

\section{不足与改进方向}
当前模型仍有简化假设，如忽略面板弹性形变、环境载荷及执行误差。后续可进一步扩展：
\begin{enumerate}
  \item 引入力学响应模型，耦合结构变形与几何拟合；
  \item 引入风载、温度场扰动，提高方案鲁棒性；
  \item 建立多目标优化模型，在接收比、能耗与安全裕度之间做综合平衡。
\end{enumerate}

\section{总结}
总体上，本文完成了从题意解析、模型建立、算法求解到结果验证的完整闭环，为主动反射面调节问题提供了系统化建模方案。
